The work flow for event selection is as follows.
Firstly the data (or Monte Carlo) collections are run over using an Athena job on the Grid.
This performs basic event reconstruction and is followed by the preselection stage.
We use release XX.XX to process the data.
The preselection software can be found in the package XX.
%The post-Athena processing code can be found in the ntuple subdirectory of the package.

In the Athena reconstruction, as a first step a JpsiFinder tool is run on muon or electron
containers, searching for di-lepton candidates with opposite sign and invariant mass
in the range [0, 6000] \mev. The \pt of the Leptons is required to be $> 2$ \gev.
%Di-lepton candidates are considered if the vertex $\chisqr$ per degree of freedom is
%$\chisqr < 30$.

\Bd candidates are then formed by BCompositeParticleBuilder [??] tool from two tracks associated
with leptons with $\pt > 500$ [??] \mev and two other opposite-signed tracks not identified as muons
or electrons (i.e. hadrons from $K^*$ candidates). These four tracks are fitted into \Bd vertex
with track mass hypotheses set for \Bd and \Bdb (kaon mass assigned to positive or negative track).
Secondary \Bd and \kstar vertices are re-fitted using TrkVKalVrtFitter [???], primary vertices
are refitted without the four candidate tracks and additional cuts are applied: 
%{\red Q) What type of fit is performed and are there any constraints used in the fit?}
\begin{itemize}
\setlength{\itemsep}{0pt}%
\setlength{\parskip}{0pt}%
 \item \Bd or \Bdb vertex $\chisqr < 205$, [??]
 \item \Bd or \Bdb mass in the range [3000, 6000] \mev, [??]
 \item \kstar or \akstar mass in the range [600, 1050] \mev. [??]
\end{itemize}

Events are saved in a linked per-event tree with one \kstar/\akstar\ candidate for each \Bd candidate:
both possible $K^\pm\pi^\pm$ assignments are considered when dealing with combinatorics and
this results in multiple candidates per event.

For the reconstruction of primary vertex (PV), the PV of the \Bd candidate is chosen to be
the one closest to the reconstructed decay vertex of the \Bd or \Bdb candidate using the three-dimensional
distance.

After the Athena reconstruction described above, a preselection is applied:
\begin{itemize}
\setlength{\itemsep}{0pt}%
\setlength{\parskip}{0pt}%
 \item opposite charged lepton pairs and hadron pairs
 \item $\Delta R_{ee} \geq 0.1$
 \item \pt(leptons)$\geq 5$~\gev
 \item $|\eta|$(leptons)$\leq 2.5$
 \item \pt(hadrons)$\geq 0.5$~\gev
 \item $|\eta|$(hadrons)$\leq 2.5$
 \item $m(ee)\leq 7$~\gev
 \item $m(K\pi ee) \in [3000,6500]$~\mev
 \item $m(K\pi) \in [690, 1110]$~\mev
 \item Track Quality requirements on all four tracks and ID requirements on electrons
% \item $|m(K\pi)-m(K^{0*})|\leq 50$~\mev
%\item \Bd$(\chisqr) \leq 2$
%\item $\tau$(\Bd)$\geq 0.2$~ps
\end{itemize}

The cut flow from this preselection for signal events ($\Bd \to \kstar ee$) and resonance events
($\Bd \to J/\psi(ee)\kstar$) are shown in Tables~\ref{tab:presel:sigele} and~\ref{tab:presel:resele}
for the electron final state.

\begin{table}[!ht] 
\caption{Number of events surviving the pre-selection for the signal non-resonant sample for the electron final state.} 
\centering 
\hspace*{-1.cm}
{\small
\begin{tabular}{l | l l l l l l l l l } 
\hline\hline
Selection &             &        & Truth- & Truth- & Relative & Relative &  &  Ratio of\\
criterion & Candidates  & Events & matched & matched & efficiency & efficiency & Purity & relative \\
          &             &        & candidates & events &          & (truth- & & efficiencies\\
          &             &        &   &   &          & matched) & & \\        
\hline
   to     &    be       &  filled      &   &   &          & & & \\        
\hline\hline
\end{tabular} 
}
\label{tab:presel:sigele}
\end{table}


\begin{table}[!ht] 
\caption{Number of events surviving the pre-selection for the reference resonant sample for the electron final state.} 
\centering 
\hspace*{-1.cm}
{\small
\begin{tabular}{l | l l l l l l l l l } 
\hline\hline
Selection &             &        & Truth- & Truth- & Relative & Relative &  &  Ratio of\\
criterion & Candidates  & Events & matched & matched & efficiency & efficiency & Purity & relative \\
          &             &        & candidates & events &          & (truth- & & efficiencies\\
          &             &        &   &   &          & matched) & & \\        
\hline
baseline & 2515085 & 217958 & 156485 & 154886 &  &  &   0.711 &  \\
$\pT(\ell_1)>5$~GeV & 2515085 & 217958 & 156485 & 154886 &  1.000 &  1.000 &  0.711 &  1.000 \\
$\pT(\ell_2)>5$~GeV & 2515085 & 217958 & 156485 & 154886 &  1.000 &  1.000 &  0.711 &  1.000 \\
$\eta(\ell_1)<2.5$ & 2515085 & 217958 & 156485 & 154886 &  1.000 &  1.000 &  0.711 &  1.000 \\
$\eta(\ell_2)<2.5$ & 2515085 & 217958 & 156485 & 154886 &  1.000 &  1.000 &  0.711 &  1.000 \\
$\pT(h_1)>1$~GeV & 913674 & 195110 & 121550 & 120286 &  0.895 &  0.777 &  0.617 &  0.868 \\
$\pT(h_2)>1$~GeV & 417361 & 153585 & 99455 & 98392 &  0.787 &  0.818 &  0.641 &  1.039 \\
$\eta(h_1)<2.5$ & 416603 & 153460 & 99312 & 98252 &  0.999 &  0.999 &  0.640 &  0.999 \\
$\eta(h_2)<2.5$ & 415912 & 153329 & 99192 & 98134 &  0.999 &  0.999 &  0.640 &  1.000 \\
$\chi^2(B)/$NDoF$<2$ & 286407 & 131015 & 88514 & 87592 &  0.854 &  0.893 &  0.669 &  1.045 \\
$\tau(B)>0.2$ & 207721 & 109944 & 77542 & 76747 &  0.839 &  0.876 &  0.698 &  1.044 \\
\hline
$|\Delta m(K\pi)_\mathrm{PDG}| < 50$~MeV & 59235 & 49955 & 29008 & 28765 &  0.454 &  0.375 &  0.576 &  0.825 \\
$m(B)\in [4.7,5.7]$~GeV & 43969 & 40179 & 26654 & 26447 &  0.804 &  0.919 &  0.658 &  1.143 \\
\hline
$|\Delta m(\pi K)_\mathrm{PDG}| < 50$~MeV & 59689 & 50276 & 29148 & 28868 &  0.457 &  0.376 &  0.574 &  0.823 \\
$m(\overline{B})\in [4.7,5.7]$~GeV & 44390 & 40545 & 26724 & 26473 &  0.806 &  0.917 &  0.653 &  1.137 \\
\hline
best $\chi^2(B)$ & 64742 & 64742 & 51785 & 51785 &  1.597 &  1.956 &  0.800 &  1.225 \\
\hline\hline
\end{tabular} 
}
\label{tab:presel:resele}
\end{table}


The cut flow from this preselection for signal events ($\Bd \to \kstar \mu\mu$) and resonance events
($\Bd \to J/\psi(\mu\mu)\kstar$) are shown in Tables~\ref{tab:presel:sigmu} and~\ref{tab:presel:resmu}
for the muon final states.

\begin{table}[!ht] 
\caption{Number of events surviving the pre-selection for the signal non-resonant sample for the muon final state.} 
\centering 
\hspace*{-1.cm}
{\small
\begin{tabular}{l | l l l l l l l l l } 
\hline\hline
Selection &             &        & Truth- & Truth- & Relative & Relative &  &  Ratio of\\
criterion & Candidates  & Events & matched & matched & efficiency & efficiency & Purity & relative \\
          &             &        & candidates & events &          & (truth- & & efficiencies\\
          &             &        &   &   &          & matched) & & \\        
\hline
   to     &    be       &  filled      &   &   &          & & & \\        
\hline\hline
\end{tabular} 
}
\label{tab:presel:sigmu}
\end{table}


\begin{table}[!ht] 
\caption{Number of events surviving the pre-selection for the reference resonant sample for the muon final state.} 
\centering 
\hspace*{-1.cm}
{\small
\begin{tabular}{l | l l l l l l l l l } 
\hline\hline
Selection &             &        & Truth- & Truth- & Relative & Relative &  &  Ratio of\\
criterion & Candidates  & Events & matched & matched & efficiency & efficiency & Purity & relative \\
          &             &        & candidates & events &          & (truth- & & efficiencies\\
          &             &        &   &   &          & matched) & & \\        
\hline
   to     &    be       &  filled      &   &   &          & & & \\        
\hline\hline
\end{tabular} 
}
\label{tab:presel:resmu}
\end{table}


The percentage of multiple candidates after the preselection is about XX\% (XX\%)
for signal (resonance) simulated events.






